\section{Experiments}

\subsection{Data and forecasting task}

We study weekly influenza hospitalization-rate forecasting using a
FluSurv-NET CSV export. The benchmark reads the raw file with
\texttt{skiprows=2}, removes non-tabular disclaimer rows with missing
\texttt{WEEK} or \texttt{YEAR.1}, sorts chronologically by year and week, and
creates a continuous time index. The primary target series is the weekly
hospitalization rate for the \emph{Entire Network} catchment with all other
metadata fields set to \emph{Overall}. We also evaluate five age-stratified
series: \texttt{0-4 yr}, \texttt{5-17 yr}, \texttt{18-49 yr},
\texttt{50-64 yr}, and \texttt{>= 65 yr}.

All epidemic models operate in normalized population units with latent
compartment states summing approximately to one, except for numerical
deviations controlled by explicit mass-conservation penalties.

\subsection{Chronological splits and metrics}

For each series, we use a chronological split with the first $60\%$ of weeks
for training, the next $20\%$ for validation, and the last $20\%$ for test.
In addition to fixed held-out test evaluation, we run rolling-origin
forecasting at horizons $1$, $2$, and $4$ weeks ahead.

Point-forecast evaluation uses:
\begin{itemize}[leftmargin=1.5em]
\item mean absolute error (MAE),
\item root mean squared error (RMSE), and
\item sMAPE.
\end{itemize}

For probabilistic forecasts, we also track negative log-likelihood, empirical
interval coverage, average interval width, and interval score. Uncertainty
calibration is handled separately at the benchmark level and does not modify
point forecasts.

\subsection{Model families}

The benchmark currently compares six model families:
\begin{enumerate}[leftmargin=1.5em]
\item \texttt{deterministic\_seir}: a seasonal discrete-time SEIR baseline with
observation model $\hat{y}_t = \rho I_t$.
\item \texttt{probabilistic\_seir}: the same latent dynamics with Student-$t$
observations and bootstrap-based uncertainty.
\item \texttt{fractional\_seir}: a fractional-order extension with shared order
$\alpha \in [0.7, 1.0]$.
\item \texttt{hospitalized\_seihr}: a hospitalization-aware baseline with
compartments $(S,E,I,H,R)$ and observation model $\hat{y}_t = \rho H_t$.
\item \texttt{delayed\_observation\_seir}: standard SEIR dynamics with delayed
observation $\hat{y}_t = \rho I_{t-d}$, where $d \in \{0,1,2,3\}$ is chosen by
validation MAE only.
\item \texttt{constrained\_structure\_discovery}: a non-LLM propose-fit-verify-
refine search over a small epidemic grammar.
\end{enumerate}

The manual \texttt{hospitalized\_seihr} baseline was corrected in the current
round to allow both direct recovery $I \rightarrow R$ and hospitalization
$I \rightarrow H \rightarrow R$, rather than forcing every infectious
individual through the hospitalized compartment.

\subsection{Observation-aware constrained discovery}

The discovery procedure searches over both latent structure and observation
semantics. The allowed compartment templates are:
\[
\texttt{SIR}, \texttt{SEIR}, \texttt{SEIRS}, \texttt{SEIHR}, \texttt{SEIAR}.
\]

The current grammar also supports:
\begin{itemize}[leftmargin=1.5em]
\item a fractional toggle,
\item observation maps \texttt{I}, \texttt{H}, \texttt{I+H}, and
\texttt{delayed\_I},
\item observation delays of $1$, $2$, or $3$ weeks for
\texttt{delayed\_I}.
\end{itemize}

Structural validity rules enforce biologically plausible infection flow,
eventual recovery, no isolated compartments, and compatibility between
observation maps and available compartments. For example, \texttt{obs=H} and
\texttt{obs=I+H} are only valid when an $H$ compartment exists, while
\texttt{obs=delayed\_I} requires an $I$ compartment.

The discovery loop is programmatic rather than language-model-driven. Starting
from \texttt{SEIR}, the procedure repeatedly:
\begin{enumerate}[leftmargin=1.5em]
\item proposes one-hop neighbors,
\item validates each candidate structure,
\item fits valid candidates,
\item scores them on validation and rolling-origin behavior,
\item retains the top beam, and
\item stops early when the search stagnates.
\end{enumerate}

To keep the search space controlled, delayed-observation neighbors are only
generated from \texttt{SIR}, \texttt{SEIR}, and \texttt{SEIRS}, while
\texttt{SEIHR} candidates prioritize \texttt{obs=I}, \texttt{obs=H}, and
\texttt{obs=I+H}. Complexity scoring includes penalties for additional free
parameters, compartments, recurrence, fractional order, observation-map
complexity, and delayed observation length.

\subsection{Seeds, robustness, and ablations}

The main robustness analysis uses five seeds:
\[
\{0, 1, 2, 42, 2024\}.
\]
For each seed, we rerun the full age-stratified benchmark and aggregate mean
performance, standard deviations, and win rates.

We also ran two ablations relevant to the current paper narrative.

\paragraph{Age-prior ablation.}
The discovery score contains an optional age-specific prior over structure
families. A single-seed age-prior ablation and a five-seed observation-aware
no-age-prior rerun both showed identical selected structures and identical
aggregate model summaries relative to the age-prior-enabled run. In the
observation-aware multi-seed setting, the ablation also leaves the dominant
observation map and dominant delay mode unchanged for every series. This
strongly suggests that the current discovery patterns are not being trivially
forced by the age prior.

\paragraph{Benchmark-level conformal calibration.}
Probabilistic interval calibration is handled as a separate benchmark-level
post-processing step, not inside model fitting. Five calibration kinds are
compared (\texttt{raw}, \texttt{scale\_calibrated},
\texttt{conformal\_absolute}, \texttt{conformal\_standardized},
\texttt{conformal\_asymmetric}), and calibration-method selection is based on
validation only. The current default winner-selection rule is the balanced
\texttt{v3} rule, which preserves nearly the same coverage gap as the most
coverage-oriented rule while reducing interval score and width.

\subsection{Implementation details}

All models use multiple random restarts and explicit penalties for negative
states and mass-conservation violations. The current age-robustness benchmark
uses six restarts for each train fit, one restart for rolling-origin refits,
and a maximum of 250 optimization iterations per restart. The discovery beam
width is five, the maximum number of rounds is twenty, and search uses both
validation MAE and rolling-origin stability criteria. The codebase is fully
tested, with plain \texttt{pytest -q} supported via a repository-level
\texttt{pytest.ini}.

\subsection{LLM-guided comparison protocol}

The repository now contains two LLM-facing layers that sit on top of the
existing non-LLM benchmark rather than replacing it.

\paragraph{LLM-V0.}
LLM-V0 is a proposal-only layer. It constructs a prompt-safe series summary,
adds a static hospitalization-semantics card, asks a proposer for JSON
structure candidates, asks a critic for annotations, applies the same hard
validator used by the non-LLM discovery pipeline, and then evaluates all
hard-valid candidates with the existing numerical executor. It therefore tests
whether proposal parsing, leakage guards, hard validation, and non-LLM
executor reuse all work end-to-end.

\paragraph{LLM-V1.}
LLM-V1 adds iterative refinement. Round 1 matches V0:
\[
\texttt{summary} \rightarrow \texttt{semantics} \rightarrow
\texttt{proposer} \rightarrow \texttt{critic} \rightarrow
\texttt{hard validator} \rightarrow \texttt{executor}.
\]
For rounds $r \geq 2$, a result analyst reads only validation- and
rolling-derived outputs from round $r-1$ and emits bounded feedback for the
next proposer call. The proposer is then restricted to small edits around the
previous best candidate, after which the critic, hard validator, and executor
run again.

\paragraph{Comparison design.}
The intended LLM comparison is three-way:
\begin{enumerate}[leftmargin=1.5em]
\item \textbf{LLM-V0}: one-shot proposal-only search;
\item \textbf{LLM-V1}: iterative validation-feedback refinement; and
\item \textbf{non-LLM constrained discovery}: the observation-aware discovery
baseline already used throughout the repository.
\end{enumerate}

This comparison is not meant to ask only whether the LLM achieves a lower test
MAE. Instead, it asks whether an agentic refinement loop improves proposal
quality, semantic alignment, and validation-selected candidate quality under
the same mechanistic constraints and executor.

\paragraph{Prompt-safe inputs and no-test leakage.}
All LLM-facing inputs are built from prompt-safe summaries. The proposer,
critic, and analyst are not allowed to see:
\begin{itemize}[leftmargin=1.5em]
\item held-out test metrics such as \texttt{test\_mae}, \texttt{test\_rmse},
or \texttt{test\_smape},
\item best-test model identifiers,
\item test-policy recommendations, and
\item test-winner files such as \texttt{benchmark\_series\_winners}.
\end{itemize}
Allowed prompt content includes only series identity, training summary,
validation summary, rolling validation summary, surveillance semantics, current
non-test discovery structure modes, and previous-round validation/rolling
feedback. Held-out test metrics are computed only after the final
validation/rolling-selected candidate is fixed.

\paragraph{Evaluation metrics.}
The LLM protocol tracks six classes of outputs:
\begin{itemize}[leftmargin=1.5em]
\item \textbf{valid proposal rate}: fraction of raw proposals that become
hard-valid candidates;
\item \textbf{candidate efficiency}: number of evaluated candidates relative to
the non-LLM reference search budget;
\item \textbf{best validation score}: the best composite
validation/rolling-oriented discovery score reached in the run;
\item \textbf{rolling validation MAE}: the mean rolling-origin validation MAE
of the selected candidate;
\item \textbf{selected-candidate test MAE}: held-out test MAE computed only
after the final validation-selected candidate is fixed; and
\item \textbf{semantic alignment}: whether proposed observation maps are
consistent with the hospitalization-rate target.
\end{itemize}
Efficiency claims are explicitly suppressed unless the LLM and non-LLM budgets
are matched.

\paragraph{Provider distinction.}
The LLM stack now supports both \texttt{provider=mock} and
\texttt{provider=openai}. The mock provider is used for engineering validation:
schema enforcement, leakage guards, round-aware artifact generation, and
executor reuse. Live-provider runs are the appropriate setting for scientific
evaluation, but only after the resulting artifacts are validated and frozen
under the same protocol. The currently frozen repository artifacts remain
mock-provider engineering runs.
